\documentclass[11pt]{article}
\usepackage{series}
\usepackage{booktabs,tabularx,array}
\usepackage{xcolor}

\newcommand{\A}{\mathcal A}
\newcommand{\C}{\mathcal C}
\newcommand{\B}{\mathcal B}
\newcommand{\Hspace}{\mathcal H}
\newcommand{\Prob}{\mathsf P}
\newcommand{\E}{\mathbb E}
\newcommand{\R}{\mathbb R}
\newcommand{\Cplx}{\mathbb C}
\newcommand{\id}{\operatorname{id}}
\newcommand{\Res}{\operatorname{Res}}
\newcommand{\States}{\mathsf S}

\title{\textbf{From Coherence-Basin Statistics to a Local Quantum-Field Interface}\\
\large Commutative Measures, the Selected Free CAR Net, and the Interacting Boundary}
\author{Peter Nero}
\date{Version 3, July 2026}

\begin{document}
\maketitle

\begin{abstract}
This paper asks a narrower and more useful question than whether quantum
field theory can be reconstructed from classical ensemble statistics.  Let
an admissible coherence-basin space carry a probability measure and bounded
basin observables.  We prove that these data have a canonical Hilbert-space
representation by multiplication operators on an \(L^2\) space.  The
representation is commutative.  Convex mixtures, covariance kernels,
Gaussian characteristic functionals, and classical coarse-graining follow
within that layer, but canonical commutation or anticommutation relations,
quantum interference, time ordering, a local net, and renormalized
interactions do not.

The correct connection to quantum field theory is an interface rather than
an emergence-by-statistics theorem.  A state on a quantum observable algebra
restricts to a probability measure on every commutative record subalgebra;
equivalently, a positive operator-valued measurement or instrument produces
classical outcome statistics.  Conversely, every state on a commutative
subalgebra has extensions to the ambient \(C^*\)-algebra, but those
extensions are generally nonunique.  Identical basin statistics can
therefore coexist with inequivalent quantum phases and correlations.

This distinction can now be applied to the current Modal Triplet Theory
(MTT) corpus.  A companion paper owns the selected free-field construction:
on the declared globally hyperbolic framed q79 representative, a twisted
massless Dirac source composes with standard CAR/AQFT machinery to give an
even local observable net with locality, covariance, the time-slice
property, and a nonempty positive Hadamard state space.  Basin and detector
statistics may be read from commuting subalgebras or instruments of that
net.  The geometry-selected nonperturbative interacting gauge-BRST
\(C^*\)-bridge, selected full-branch state, renormalization-group matching,
and observable uncertainty packet remain open.  The result is a precise
classical-to-quantum interface and a closed selected free-QFT source, not a
derivation of noncommutative QFT from probability alone.
\end{abstract}

\section*{Version 3 revision note}

\begin{description}[style=nextline]
\item[Supersedes.]
Version 2, DOI \url{https://doi.org/10.5281/zenodo.18322089}.
\item[Reason.]
Version 2 claimed that coherence-basin measures generated Hilbert-space
quantum states, noncommuting fields, propagators, path integrals,
renormalization, and Feynman rules.  Its GNS argument began with a
commutative algebra and then inserted noncommuting operators without a
construction.  It also conflated classical covariance with several
different quantum and causal kernels.
\item[Resolution.]
Version 3 proves the exact commutative representation supplied by basin
statistics, gives a finite-dimensional nonuniqueness witness, and replaces
the unsupported reconstruction with a typed state-restriction and
state-extension interface.  It imports, rather than duplicates, the
separately owned selected free-CAR theorem and records the current
interacting boundary.
\item[Retained content.]
Coherence basins can still provide useful configuration variables,
preparation distributions, detector records, classical correlations, and
coarse-grained summaries.  Quantum fields need not be interpreted as
classical substances.  What is withdrawn is the claim that ordinary
probability theory alone creates their noncommutative algebra.
\item[Open boundary.]
The selected free q79 CAR net is closed at its declared tier.  A unique
state-selection rule and a geometry-selected nonperturbative interacting
gauge-BRST \(C^*\)-completion remain open, as do full RG matching and
precision-observable comparison.
\end{description}

\tableofcontents

\section{The question after the current MTT advances}

There are two very different claims that can be made about a statistical
description.
\begin{enumerate}
\item A collection of physical configurations and a probability measure
can be represented by functions and multiplication operators on a Hilbert
space.
\item Those statistical data determine a noncommutative quantum theory with
field equations, causal propagators, local observable algebras, states, and
interactions.
\end{enumerate}
The first claim is a standard and exact representation theorem.  The second
does not follow from the first.

This distinction matters especially in MTT.  The earlier version of this
paper tried to make basin statistics do all of the work.  The current
research program has since supplied a more concrete route.  A selected
geometric source gives a Dirac-type operator; established CAR and algebraic
QFT constructions then give the free quantum observable net.  Statistics
enter when a state and a physical record channel are applied to that net.

The revised chain is therefore
\[
\begin{gathered}
\text{selected geometry and bundle}
\longrightarrow
\text{Green-hyperbolic field operator}\\
\longrightarrow
\text{CAR observable net},\\
\text{state and record channel}
\longrightarrow
\text{classical basin or detector law}.
\end{gathered}
\]
The lower line is not a lesser physical event.  Preparation, interaction
with an apparatus, amplification, and recording are ordinary physical
processes.  The point is only that their classical output law does not,
when read backwards, uniquely determine the quantum algebra that produced
it.

\section{The classical basin layer}

\subsection{Typed statistical data}

Let \((\B,\Sigma,\mu)\) be a probability space whose points represent
admissible coherence-basin configurations at a declared resolution.  A
bounded classical basin observable is an element
\[
 f\in L^\infty(\B,\Sigma,\mu).
\]
Its expectation is
\[
 \E_\mu[f]=\int_{\B}f(b)\,d\mu(b).
\]
The observable algebra
\[
 \C_\mu:=L^\infty(\B,\Sigma,\mu)
\]
is a commutative unital \(C^*\)-algebra up to the usual identification
almost everywhere.  If topological basin data are needed, one may instead
start from \(C(K)\) for a compact Hausdorff configuration space \(K\).

These definitions are intentionally modest.  They describe a probability
law over configurations and the functions that can be averaged.  They do
not assume that a basin is a particle, that \(\mu\) is fundamental, or that
the basin variables are complete.

\subsection{The canonical Hilbert representation}

\begin{theorem}[Commutative basin representation]
\label{thm:commutative}
Let \((\B,\Sigma,\mu)\) be a probability space.  On
\(\Hspace_\mu=L^2(\B,\mu)\), define
\[
 (M_f\psi)(b):=f(b)\psi(b),
 \qquad f\in L^\infty(\B,\mu).
\]
Then \(f\mapsto M_f\) is a faithful unital \(*\)-representation of
\(\C_\mu\), the vector \(\mathbf 1\in\Hspace_\mu\) represents the
expectation state,
\[
 \E_\mu[f]=\langle\mathbf 1,M_f\mathbf 1\rangle,
\]
and
\[
 [M_f,M_g]=0
 \qquad\text{for every }f,g\in\C_\mu.
\]
\end{theorem}

\begin{proof}
Multiplication by an essentially bounded function is a bounded operator on
\(L^2\), with \(M_f^*=M_{\overline f}\),
\(M_fM_g=M_{fg}\), and \(M_1=I\).  If \(M_f=0\), then \(f=0\) almost
everywhere, so the representation is faithful.  The expectation identity
is immediate from the \(L^2\) inner product, and multiplication operators
commute because \(fg=gf\).
\end{proof}

This is the relevant GNS representation in concrete form.  It explains why
Hilbert spaces occur naturally in probability theory.  It does not turn a
commutative algebra into a noncommutative one.

\begin{corollary}[Probability does not supply CCR or CAR]
\label{cor:no-car}
The data of Theorem~\ref{thm:commutative} do not by themselves produce
operators satisfying nonzero canonical commutation relations or canonical
anticommutation relations.  Any such operators require additional algebraic
and dynamical input.
\end{corollary}

\begin{proof}
Every represented basin observable lies in a commutative algebra.  In
particular, all commutators vanish.  A CAR algebra also contains odd
generators with prescribed anticommutators and a grading.  Neither structure
is specified by the probability space.
\end{proof}

\subsection{Convexity is not coherent superposition}
\leavevmode\par

If \(\mu_1\) and \(\mu_2\) are probability measures, then
\[
 \mu_t=t\mu_1+(1-t)\mu_2,\qquad 0\leq t\leq1,
\]
is a statistical mixture.  Expectations depend affinely on \(t\).
Quantum density matrices are also convex, but a coherent vector
\[
 \psi=\alpha\psi_1+\beta\psi_2
\]
contains relative phase information that is not determined by the
probabilities \(|\alpha|^2\) and \(|\beta|^2\).  Off-diagonal matrix
elements can change interference observables while every diagonal outcome
probability remains fixed.

Classical overlap of two probability densities can therefore model
ordinary uncertainty or common support.  It is not, without a phase-bearing
quantum algebra and state, a derivation of interference.

\section{Why the reverse reconstruction is nonunique}

\subsection{A two-level witness}

The obstruction is visible in the smallest noncommutative example.  In
\(\A=M_2(\Cplx)\), let
\[
 \C=\left\{
 \begin{pmatrix}a&0\\0&b\end{pmatrix}:a,b\in\Cplx
 \right\}
 \cong C(\{0,1\}).
\]
For \(t\in[-1/2,1/2]\), define
\[
 \rho_t=
 \begin{pmatrix}
 1/2&t\\
 t&1/2
 \end{pmatrix}.
\]
Each \(\rho_t\) is a density matrix and gives the same probability
\((1/2,1/2)\) on \(\C\).  Yet
\[
 \tr(\rho_t\sigma_x)=2t
\]
depends on \(t\).  A fixed classical outcome law therefore leaves a whole
interval of quantum coherences undetermined.

\begin{proposition}[Classical-law nonselection]
\label{prop:nonselection}
There is no rule depending only on a classical probability space
\((\B,\Sigma,\mu)\) that uniquely selects an ambient noncommutative
\(C^*\)-algebra and state up to physical equivalence.
\end{proposition}

\begin{proof}
The two-level construction already gives distinct states on one fixed
noncommutative algebra with the same restriction to its classical diagonal
subalgebra.  Tensoring with additional noncommutative systems gives further
ambient algebras with the same classical restriction.  Hence the classical
law does not determine either the off-diagonal state data or the full
ambient algebra.
\end{proof}

The proposition does not say that classical records are uninformative.
Tomography can recover a quantum state when a sufficiently rich family of
noncommuting preparations and effects is already defined.  That procedure
uses a specified quantum model and an informationally complete
measurement; it does not infer the quantum model from one classical law.

\subsection{What a valid reconstruction would need}

To reconstruct a QFT rather than a classical random field, one must provide
at least:
\begin{enumerate}
\item a spacetime category or a fixed causal spacetime;
\item a field equation or local action with a declared domain;
\item a symplectic or Hermitian solution-space structure;
\item a CCR, CAR, or other observable algebra;
\item local covariance, causality, and time-slice maps;
\item a positive state satisfying the relevant spectrum or Hadamard
condition;
\item for gauge fields, the BRST/BV quotient and physical observables;
\item for interactions, a renormalized product or controlled continuum
construction.
\end{enumerate}
Probability distributions may be outputs of this structure.  They cannot
replace these eight items.  Standard operator-algebraic and locally
covariant formulations make these obligations explicit
\cite{bratteli_robinson,haag,fewster_rejzner}.

\section{The exact classical-to-quantum interface}

\subsection{Restriction of a quantum state}

Let \(\A\) be a unital \(C^*\)-algebra and let
\(\C\subseteq\A\) be a unital commutative \(C^*\)-subalgebra.  By the
Gelfand representation, \(\C\cong C(K)\) for a compact Hausdorff spectrum
\(K\).

\begin{theorem}[State restriction and extension interface]
\label{thm:interface}
The restriction map
\[
 \Res_{\C}:\States(\A)\longrightarrow\States(\C),
 \qquad
 \omega\longmapsto\omega|_{\C},
\]
is affine.  Every restricted state has a unique probability measure
\(\mu_\omega\) on \(K\) satisfying
\[
 \omega(c)=\int_K\widehat c(k)\,d\mu_\omega(k),
 \qquad c\in\C.
\]
Conversely, every state on \(\C\) has at least one state extension to
\(\A\).  The extension need not be unique.
\end{theorem}

\begin{proof}
Affineness is immediate.  The Riesz--Markov representation theorem gives
the unique probability measure associated with the positive normalized
functional \(\omega|_{\C}\).  The positive Hahn--Banach state-extension
theorem extends any state on a unital \(C^*\)-subalgebra to a state on
\(\A\) \cite{bratteli_robinson}.  Nonuniqueness is witnessed by the family
\(\rho_t\) above.
\end{proof}

This theorem is the corrected mathematical meaning of ``basin statistics
as a shadow of quantum theory.''  The arrow from a quantum state to a
classical law is exact once the record subalgebra has been chosen.  The
reverse arrow exists as an extension problem but is generally multivalued.

\subsection{POVMs and instruments}

A detector need not correspond to a sharp commutative subalgebra.  A
positive operator-valued measure \(E\) on an outcome space \(K\) assigns
effects \(E(\Delta)\in\A\) to measurable outcome sets.  In state \(\omega\),
\[
 \mu_\omega(\Delta)=\omega(E(\Delta))
\]
is a probability measure.  A quantum instrument additionally specifies
the state update associated with each outcome.  Both are physical maps
implemented by an interaction, amplification, and record process.

Nothing in this description makes measurement metaphysically privileged.
It is a particular coupling between systems that produces a durable
classical record.  What matters mathematically is the chosen effect or
instrument, not the presence of an observer.

\subsection{Local compatibility}

For a local QFT net \(O\mapsto\A(O)\), choose compatible record
subalgebras \(\C(O)\subseteq\A(O)\) or local instruments.  A net state
\(\omega\) then supplies compatible local probability laws
\[
 \mu_{\omega,O}:=\omega|_{\C(O)}.
\]
If \(O_1\subseteq O_2\) and the inclusions of record algebras commute with
the net inclusions, the restricted laws have the corresponding consistency.

Independent state extensions region by region do not automatically glue
to one global net state.  Compatibility, causality, and the selected state
condition remain global obligations.

\section{Four kernels that must remain distinct}

The earlier paper used the word ``propagator'' for a classical covariance.
Several kernels can have similar formulas while carrying different
structures.

\begin{center}
\small
\begin{tabularx}{\textwidth}{>{\raggedright\arraybackslash\bfseries}p{0.23\textwidth}>{\raggedright\arraybackslash}p{0.30\textwidth}>{\raggedright\arraybackslash}X}
\toprule
Kernel & Definition or source & What it controls \\
\midrule
Classical covariance &
\(\E_\mu[(\rho(x)-\bar\rho(x))(\rho(y)-\bar\rho(y))]\) &
Second moments of a classical random field. \\
Advanced/retarded Green kernel &
Fundamental solution \(G^\pm\) of a Green-hyperbolic operator &
Causal response and support in \(J^\pm\). \\
Quantum two-point function &
\(W_\omega(x,y)=\omega(\phi(x)\phi(y))\) in a state \(\omega\) &
State correlations, positivity, field equation, and singularity structure. \\
Feynman kernel &
Time-ordered two-point distribution or perturbative inverse &
Perturbative contractions subject to a prescription and state. \\
\bottomrule
\end{tabularx}
\end{center}

A classical Gaussian covariance can resemble a Euclidean free-field
two-point function.  A QFT reconstruction from Euclidean data nevertheless
requires additional conditions such as reflection positivity, covariance,
regularity, and clustering.  In Lorentzian AQFT, the commutator or
anticommutator, microlocal spectrum condition, and positivity are separate
requirements \cite{osterwalder_schrader1,osterwalder_schrader2,
radzikowski,bfv}.

\begin{proposition}[Covariance is not a propagator certificate]
A symmetric positive classical covariance kernel does not by itself
determine retarded and advanced Green operators, a quantum commutator or
anticommutator, a time-ordered product, or a local observable net.
\end{proposition}

\begin{proof}
The classical covariance contains no causal orientation, field operator,
grading, or local-algebra assignment.  Distinct quantum states and even
distinct field equations can have the same covariance on a restricted
commuting record family.  The missing structures cannot be recovered from
positivity and symmetry alone.
\end{proof}

\section{Characteristic functionals and path integrals}

For a classical random distribution \(\rho\), the functional
\[
 Z_{\mathrm{cl}}[J]
 =\E_\mu\!\left[e^{i\rho(J)}\right]
\]
is its characteristic functional.  Functional derivatives generate
classical moments when the derivatives exist.  This is useful and exact.

A Lorentzian QFT path integral is not generally a countably additive
probability measure with density \(e^{iS}\).  In perturbation theory it is
often a compact notation for time-ordered distributions, Wick expansion,
gauge fixing, renormalization, and boundary conditions.  Euclidean
functional measures can be rigorous in selected models, but passage to a
Lorentzian QFT uses a reconstruction theorem and its hypotheses.

\begin{proposition}[Generating-functional boundary]
The existence of \(Z_{\mathrm{cl}}\) supplies classical moments.  It
becomes evidence for a quantum field theory only after an independent
reconstruction verifies the required algebra, causal or Euclidean
structure, positivity, and state conditions.
\end{proposition}

This is not a defect of basin statistics.  A classical random field may be
the right effective model for a detector signal, stochastic environment,
hydrodynamic variable, or semiclassical noise source.  It should simply be
named correctly.

\section{Coarse-graining, renormalization, and diagrams}

\subsection{Three different reductions}

It is useful to separate:
\begin{enumerate}
\item \textbf{statistical marginalization}, which integrates out random
variables in a probability measure;
\item \textbf{projection or truncation}, which retains a selected effective
subspace or finite set of observables;
\item \textbf{renormalization-group transport}, which changes scale while
preserving a declared class of observables through running couplings,
counterterms, or effective actions.
\end{enumerate}
They may be combined in one model, but they are not synonyms.

A coherence capacity can control a domain, truncation error, spectral gap,
or first-exit margin.  It does not alone calculate beta functions.  A
Wilsonian or perturbative RG flow requires a specified action or local
algebra, scale prescription, regulator or subtraction scheme, and matching
conditions.

\subsection{Where Feynman diagrams come from}

Feynman diagrams encode the combinatorics of a perturbative expansion once
the free propagator, interaction vertices, grading, time ordering, and
renormalization prescription have been specified.  Basin correlations can
be represented diagrammatically too, as in cumulant or cluster expansions.
The visual similarity does not establish equality of the theories.

Within the current MTT Standard-Model branch, the perturbative observable
functor uses standard BRST/Faddeev--Popov quantization.  That layer is an
imported standard-QFT construction composed with the selected carrier, not
a derivation of quantization from coherence-basin counting.

\section{The selected MTT quantum-field source}

\subsection{The theorem owned by the companion paper}

The paper \emph{Modal Triplet Theory and Quantum Field Theory on Curved
Spacetime: A Selected Free CAR Net and the Interacting Reconstruction
Boundary} owns the current selected free-QFT theorem
\cite{nero_mtt_qft}.  Its declared source is a globally hyperbolic framed
q79 representative with a global coframe and parallel rank-six carrier.
Those data define a zero-order-free twisted massless Dirac operator.  Under
the standard Green-hyperbolic and CAR/AQFT theorems, the construction gives:
\begin{itemize}
\item an even CAR observable net;
\item locality, covariance, and the time-slice property;
\item a nonempty positive Hadamard state space;
\item an exact finite coherent/complement component map.
\end{itemize}
The independent analytic and algebraic ingredients are standard results for
Green-hyperbolic equations, Dirac fields, local covariance, and Hadamard
renormalization
\cite{bar_ginoux_pfaeffle,dimock,bfv,radzikowski,hollands_wald}.

This paper does not reprove or rename that result.  It explains how
coherence-basin and detector statistics can be attached to it through
Theorem~\ref{thm:interface}.  A selected state \(\omega\) and a record
subalgebra or instrument produce the corresponding probability law.

\subsection{What is closed and what is not}

\begin{center}
\small
\begin{tabularx}{\textwidth}{>{\raggedright\arraybackslash\bfseries}p{0.30\textwidth}>{\raggedright\arraybackslash}p{0.18\textwidth}>{\raggedright\arraybackslash}X}
\toprule
Object & Current tier & Meaning for this paper \\
\midrule
Finite selected quantum model & Closed &
The finite-symbol Cauchy Hilbert/state/observable system and its declared
comparison maps are available. \\
Selected local QFT source & Closed &
The q79 twisted-Dirac even-CAR net supplies the quantum algebra used by the
interface. \\
Basin probability representation & Exact &
Theorem~\ref{thm:commutative}; commutative and not a quantum source. \\
Quantum-to-record restriction & Exact &
Theorem~\ref{thm:interface}; the reverse extension is nonunique. \\
Perturbative SM observable functor & Conditional, imported &
Standard BRST/Faddeev--Popov quantization is used rather than derived from
MTT. \\
Finite-domain quantization results & Conditional, finite domain &
Useful controlled results exist, with explicit continuum and gauge
obligations. \\
Nonperturbative interacting QFT & Open &
Requires a geometry-selected gauge-BRST \(C^*\)-bridge or a controlled
regulator/continuum limit. \\
Unique full-branch state & Open &
Nonempty Hadamard state space does not select one physical state. \\
\bottomrule
\end{tabularx}
\end{center}

\subsection{The interacting boundary}

The current interacting program contains a classical BV master action, a
gauge-fixed Green-hyperbolic equicausal algebra, a formal all-orders
anomaly-free quantum master equation, a formal physical-state functor, and a
literal free \(C^*\)-reference net at zero coupling.  These are meaningful
structural advances.

They do not select a unique fixed-nonzero-coupling \(C^*\)-completion.  The
current flat-deformation result shows why a formal power-series jet is
insufficient: distinct nonperturbative theories can share the same formal
expansion.  The remaining exit is a same-source nonperturbative gauge-BRST
bridge or regulator/continuum limit, followed by state selection,
renormalization-group matching, uncertainty control, and observable
comparison.

\section{Physical interpretation without overclaiming}

\subsection{What basin variables can mean}

A coherence basin may serve as:
\begin{itemize}
\item a classical label for a stable projected configuration;
\item a preparation class;
\item a detector-record class;
\item a coarse-grained feature of a quantum state;
\item a semiclassical or stochastic effective variable.
\end{itemize}
These roles are not automatically identical.  Each application should name
the source map and the observable domain.

In particular, a particle in QFT is not generally a localized classical
basin.  Particle interpretations depend on representations, asymptotic
structure, or detector models, especially in curved spacetime.  A basin may
encode a stable effective particle-like regime while the local quantum
field remains the algebraic source of correlations.

\subsection{Measurement as a physical interface}

A measurement is an interaction that correlates a system with an apparatus
and produces a record.  In the algebraic description, an effect or
instrument maps a state to outcome probabilities and conditional states.
In the basin description, the same output may be represented by a
probability measure over stable records.

This does not make measurement the act that creates reality.  It identifies
the precise physical channel at which noncommutative state information is
converted into a commutative record law.  Whether MTT supplies a deeper
same-source capture dynamics is a separate theorem.

\subsection{Ontology remains optional}

The mathematics is compatible with several ontological readings.  One may
treat quantum fields as fundamental, as effective local observables, or as
linear shadows of deeper closure dynamics.  The interface theorem does not
decide among those readings.  It does decide a mathematical point: a
classical ensemble does not uniquely determine the noncommutative theory
above it.

\section{Validation program}

A proposed basin-to-QFT model should provide the following certificate.
\begin{enumerate}
\item \textbf{Basin domain:} the measurable or topological configuration
space and the meaning of its points.
\item \textbf{Probability source:} a preparation law, state restriction,
instrument, or dynamical derivation of \(\mu\).
\item \textbf{Quantum source:} the field operator, solution structure, and
CCR/CAR or local algebra.
\item \textbf{Interface map:} the commuting subalgebra, POVM, instrument, or
channel producing the basin records.
\item \textbf{State condition:} positivity and the appropriate spectrum,
Hadamard, KMS, or asymptotic condition.
\item \textbf{Causal control:} locality, covariance, and time-slice
behavior.
\item \textbf{Interaction control:} BV/BRST identities, renormalization,
continuum or nonperturbative construction, and uncertainty estimates.
\item \textbf{Observable comparison:} a prediction not used to select the
source data.
\end{enumerate}

The model is falsified at its declared tier if the record law cannot be
obtained from the selected state and interface, if the claimed local net
violates causality or positivity, if a stated state is not Hadamard on the
declared spacetime, or if the controlled observable comparison fails.

\section{Discussion}

\subsection{What the correction gains}

The earlier claim was rhetorically broad but mathematically incomplete:
\[
 \text{basin measure}
 \;\not\Longrightarrow\;
 \text{noncommutative QFT}.
\]
The corrected picture has more structure:
\[
\begin{gathered}
\text{q79 geometry}
\longrightarrow
\text{Dirac source}
\longrightarrow
\text{free CAR net},\\
\text{state + physical record channel}
\longrightarrow
\text{basin statistics}.
\end{gathered}
\]
Every arrow now has a type.  The free quantum source is not inferred from a
classical covariance.  The statistical layer is not discarded; it is
placed where it can be exact.

\subsection{Why this is still an MTT result}

The correction does not reduce the paper to a generic textbook observation.
The generic theorems identify the only valid interface.  Current MTT
research supplies a selected source on one side of that interface: the
q79 twisted-Dirac free CAR net.  MTT also supplies finite projection and
apparatus maps on declared domains.  The typed upper/coherent/effective
distinction and the exact limits of projection claims are stated in the
revised foundation and projection papers
\cite{nero_foundations,nero_projection}.  The combined statement is
therefore specific:
\begin{quote}
The selected free MTT quantum-field source can emit ordinary basin and
detector statistics through compatible record maps, but those statistics do
not independently reconstruct the source.
\end{quote}

That statement is stronger than the old statistical analogy because it
names the quantum object being represented.  It is narrower than a full QFT
derivation because the interacting completion and unique state selection
remain open.

\subsection{The best next theorem}

The next major step is not another Hilbert representation of classical
probability.  It is a same-source construction that carries the selected
q79 geometry through:
\[
\begin{gathered}
 \text{interacting BV/BRST data}
 \longrightarrow
 \text{physical gauge-invariant }C^*\text{-net}\\
 \longrightarrow
 \text{selected state and record instrument}.
\end{gathered}
\]
with a regulator or renormalized pushforward that preserves the quantum
master equation, locality, positivity, and the observable comparison in the
continuum limit.

\section{Conclusion}

Coherence-basin statistics possess a natural Hilbert-space representation,
but it is a commutative one.  They support expectations, mixtures,
covariances, characteristic functionals, and classical coarse-graining.
They do not by themselves supply quantum phases, CCR/CAR relations,
time-ordered propagators, local nets, gauge reduction, or renormalized
interactions.

The correct bridge is an interface.  A quantum state restricts to an exact
classical probability law on a commuting record algebra or through an
instrument.  Classical states extend back to the ambient algebra, but
generally in many inequivalent ways.  The reverse reconstruction is
therefore nonunique.

For current MTT, this interface attaches to a real selected result: the
free q79 twisted-Dirac even-CAR net with locality, covariance, time-slice
behavior, and positive Hadamard states.  The interacting nonperturbative
gauge-BRST \(C^*\)-completion, selected full-branch state, RG matching, and
precision comparison remain the frontier.  This is the defensible
achievement of the revised paper: not QFT from statistics alone, but a
typed and testable connection between a selected quantum-field source and
its classical coherence-basin records.

\bibliographystyle{plain}
\bibliography{main}

\end{document}
